
\documentclass{article}
\usepackage{colortbl}
\usepackage{makecell}
\usepackage{multirow}
\usepackage{supertabular}

\begin{document}

\newcounter{utterance}

\centering \large Interaction Transcript for game `chronicle', experiment `political\_easy', episode 4 with Bentel rockers\_\_qwen3.5{-}4b\_\_Bentel\_iter.
\vspace{24pt}

{ \footnotesize  \setcounter{utterance}{1}
\setlength{\tabcolsep}{0pt}
\begin{supertabular}{c@{$\;$}|p{.15\linewidth}@{}p{.15\linewidth}p{.15\linewidth}p{.15\linewidth}p{.15\linewidth}p{.15\linewidth}}
    \# & $\;$A & \multicolumn{4}{c}{Game Master} & $\;\:$B\\
    \hline

    \theutterance \stepcounter{utterance}  
    & & \multicolumn{4}{p{0.6\linewidth}}{
        \cellcolor[rgb]{0.9,0.9,0.9}{
            \makecell[{{p{\linewidth}}}]{
                \texttt{\tiny{[P1$\langle$GM]}}
                \texttt{You are playing a historical deduction game as the NARRATOR.} \\
\\ 
\texttt{You have been given the name of a historical event, discovery, or milestone. Your goal is to help the DETECTIVE figure out which event you are describing — but you must describe it indirectly, without ever naming it.} \\
\\ 
\texttt{Rules:} \\
\texttt{1. Each round, write exactly ONE sentence describing the event. Add NEW information each round — do not repeat what you wrote before.} \\
\texttt{2. You must NEVER use any of the following forbidden keywords (these would make the event too easy to identify): allied, american, dachau, trials} \\
\texttt{3. You must NEVER state the event name directly or any close variant of it.} \\
\texttt{4. You may freely describe causes, consequences, the broader context, the sequence of events, and the significance — but only using general, non{-}identifying language.} \\
\\ 
\texttt{Format your response EXACTLY as follows (required):} \\
\texttt{PARAGRAPH: <exactly one sentence describing the event without forbidden keywords>} \\
\\ 
\texttt{The event you must describe (without naming it) is: Dachau trials} \\
            }
        }
    }
    & \\ \\

    \theutterance \stepcounter{utterance}  
    & \multicolumn{4}{p{0.6\linewidth}}{
        \cellcolor[rgb]{0.9,0.9,0.9}{
            \makecell[{{p{\linewidth}}}]{
                \texttt{\tiny{[P1$\rangle$GM]}}
                \texttt{PARAGRAPH: This series of legal proceedings examined the actions of individuals who had overseen the systematic extermination of people in a remote facility near a German border town.} \\
            }
        }
    }
    & & \\ \\

    \theutterance \stepcounter{utterance}  
    & & \multicolumn{4}{p{0.6\linewidth}}{
        \cellcolor[rgb]{0.9,0.9,0.9}{
            \makecell[{{p{\linewidth}}}]{
                \texttt{\tiny{[GM$\rangle$P2]}}
                \texttt{You are playing a historical deduction game as the DETECTIVE.} \\
\\ 
\texttt{A Narrator will describe a historical event, discovery, or milestone to you over several rounds. Each round you receive a new sentence of description that builds on what came before. Your job is to figure out which specific event is being described.} \\
\\ 
\texttt{After EACH new sentence, respond briefly. Keep your entire response under 100 words. Use this exact format (both sections are required):} \\
\\ 
\texttt{ANALYSIS: <2–3 short sentences: what does the growing description suggest? What type of event is this? What are the key features you are noticing?>} \\
\\ 
\texttt{GUESS: <your best guess for the name of the event being described — be specific, e.g. "The French Revolution" not just "a revolution">} \\
\\ 
\texttt{Rules:} \\
\texttt{{-} You may refine your GUESS each round as more information arrives.} \\
\texttt{{-} Your GUESS should name the specific event, not just a broad category.} \\
\texttt{{-} You have up to 10 rounds total.} \\
\texttt{{-} Be concise — quality of reasoning matters more than length.} \\
\\ 
\texttt{The first sentence from the Narrator is coming now.} \\
\\ 
\\ 
\texttt{{[}Sentence 1{]}} \\
\texttt{This series of legal proceedings examined the actions of individuals who had overseen the systematic extermination of people in a remote facility near a German border town.} \\
            }
        }
    }
    & \\ \\

    \theutterance \stepcounter{utterance}  
    & & & \multicolumn{4}{p{0.6\linewidth}}{
        \cellcolor[rgb]{0.9,0.9,0.9}{
            \makecell[{{p{\linewidth}}}]{
                \texttt{\tiny{[GM$\langle$P2]}}
                \texttt{ANALYSIS: The description points to a major historical trial involving high{-}ranking officials responsible for a genocide at a remote camp near a German border. The key features are "legal proceedings," "systematic extermination," and the specific geographic clue. This suggests the Nuremberg Trials or a specific chamber trial within that series.} \\
\\ 
\texttt{GUESS: The Nuremberg Trials} \\
            }
        }
    }
     \\ \\

    \theutterance \stepcounter{utterance}  
    & & \multicolumn{4}{p{0.6\linewidth}}{
        \cellcolor[rgb]{0.9,0.9,0.9}{
            \makecell[{{p{\linewidth}}}]{
                \texttt{\tiny{[P1$\langle$GM]}}
                \texttt{Round 1 complete. The Detective's current guess: "The Nuremberg Trials". Please write a new paragraph with fresh information.} \\
            }
        }
    }
    & \\ \\

    \theutterance \stepcounter{utterance}  
    & \multicolumn{4}{p{0.6\linewidth}}{
        \cellcolor[rgb]{0.9,0.9,0.9}{
            \makecell[{{p{\linewidth}}}]{
                \texttt{\tiny{[P1$\rangle$GM]}}
                \texttt{PARAGRAPH: These hearings specifically addressed the crimes committed at a concentration camp that was the first to be established by the Nazi regime in 1933.} \\
            }
        }
    }
    & & \\ \\

    \theutterance \stepcounter{utterance}  
    & & \multicolumn{4}{p{0.6\linewidth}}{
        \cellcolor[rgb]{0.9,0.9,0.9}{
            \makecell[{{p{\linewidth}}}]{
                \texttt{\tiny{[GM$\rangle$P2]}}
                \texttt{{[}Sentence 2{]}} \\
\texttt{These hearings specifically addressed the crimes committed at a concentration camp that was the first to be established by the Nazi regime in 1933.} \\
            }
        }
    }
    & \\ \\

    \theutterance \stepcounter{utterance}  
    & & & \multicolumn{4}{p{0.6\linewidth}}{
        \cellcolor[rgb]{0.9,0.9,0.9}{
            \makecell[{{p{\linewidth}}}]{
                \texttt{\tiny{[GM$\langle$P2]}}
                \texttt{ANALYSIS: The description narrows the focus to a specific camp established by the Nazis in 1933, which was the first of its kind. This geographic and chronological detail points directly to Dachau, distinguishing it from other camps like Auschwitz or Treblinka.} \\
\\ 
\texttt{GUESS: The Dachau Trial} \\
            }
        }
    }
     \\ \\

    \theutterance \stepcounter{utterance}  
    & & & \multicolumn{2}{p{0.3\linewidth}}{
        \cellcolor[rgb]{0.9,0.9,0.9}{
            \makecell[{{p{\linewidth}}}]{
                \texttt{\tiny{[GM$|$GM]}}
                \texttt{round 2} \\
            }
        }
    }
    & & \\ \\

\end{supertabular}
}

\end{document}
